\chapter[Photon Eigenbranch and Vacuum Relaxation]
{Dispersion-Suppressed Photon Eigenbranch and Coherent-Action-Preserving Vacuum Relaxation}
\label{app:photon_dispersion_relaxation}

\noindent\textbf{Abstract.}
This appendix develops a finite-wave-number propagation closure for the transverse optical branch of the active Face--Centered Cubic vacuum and separates conservative lattice dispersion from nonconservative vacuum relaxation. It inherits the objective FCC contact scaffold, primitive-cell inertia, and long-wavelength transverse normalization established in \hyperref[app:fcc_microphysics]{Appendix~\ref*{app:fcc_microphysics}}; the coherent Condensate constraint and localized impedance hierarchy established in \hyperref[app:condensate_impedance]{Appendix~\ref*{app:condensate_impedance}}; and the gravitational endpoint-transfer law established in \hyperref[app:gravitational_endpoint_transfer]{Appendix~\ref*{app:gravitational_endpoint_transfer}}.

A positive reduced lattice action is constructed on the divergence-free transverse constraint manifold by adding the smallest centered-second-difference term that removes the leading quadratic correction to the photon group velocity. Within the selected two-term stencil, the required coefficient is \(1/12\). The resulting branch preserves
\(c_{\mathrm{P}}=c_{\mathrm{E}}(1-\alpha_{\mathrm{G}})^2\), remains gapless and spectrally nonnegative, and shifts the leading group-speed correction from order \(k^2d_{\mathrm{nn}}^2\) to order \(k^4d_{\mathrm{nn}}^4\).

The nonconservative construction then distinguishes ordinary amplitude attenuation from coherent frequency transport. Frequency evolution requires both modal-energy transfer and preservation of coherent modal action, together with a direction-preserving transport law in wave-vector space. The finite-wave-number loading inheritance, inverse-frequency loss tangent, and coherent-action-preserving transport law are explicit constitutive closures whose recovery from the complete FCC--Condensate parent dynamics remains open.

\section{Purpose, State Domain, and Logical Status}
\label{sec:photon_relaxation_scope}

The active TFT vacuum is a dual-phase mechanical medium composed of the solid FCC Soft Floor and the connected interstitial Condensate. The gravity-neutral FCC contact action, primitive-cell inertia, affine acoustic tensor, and baseline directional branch are developed in \hyperref[app:fcc_microphysics]{Appendix~\ref*{app:fcc_microphysics}}. That construction does not complete the full finite-wave-number spectrum or its rotational and internal-phase branches.

The conservative coherent sector established in \hyperref[app:condensate_impedance]{Appendix~\ref*{app:condensate_impedance}} satisfies
\begin{equation}
    \mathbf{v}_{\mathrm{rel}}
    =
    \mathbf{v}_{\mathrm{IC}}
    -
    \dot{\mathbf{u}}
    =
    \mathbf{0},
    \qquad
    \nabla\cdot\mathbf{u}
    =
    0.
    \label{eq:photon_relaxation_coherent_limit}
\end{equation}
The Condensate is therefore entrained with the FCC scaffold at the retained carrier order and supplies the pressure response enforcing the transverse constraint. Relative flow, dissipative impedance, and irreversible reservoir transfer belong to a separate nonconservative sector and are not duplicated as carrier inertia.

The carrier action developed below is restricted to the divergence-free coherent subspace. Its admissible fields and variations obey
\begin{equation}
    \nabla\cdot\mathbf{u}
    =
    0,
    \qquad
    \nabla\cdot\delta\mathbf{u}
    =
    0.
    \label{eq:photon_relaxation_transverse_variation}
\end{equation}
Its Euler--Lagrange equation therefore governs only the selected transverse sector; it is not an unconstrained three-polarization lattice theory.

The analysis is restricted to a locally uniform active FCC background satisfying
\begin{equation}
    0
    \leq
    \alpha_{\mathrm{G}}
    <
    1,
    \qquad
    \dot{\alpha}_{\mathrm{G}}
    \approx
    0.
    \label{eq:photon_relaxation_state_domain}
\end{equation}
The gravitational loading is quasistatic during the local optical disturbance. The corresponding long-wavelength propagation capacity inherited from \hyperref[app:gravitational_endpoint_transfer]{Appendix~\ref*{app:gravitational_endpoint_transfer}} is
\begin{equation}
    c_{\mathrm{P}}
    =
    c_{\mathrm{E}}
    \left(
        1-\alpha_{\mathrm{G}}
    \right)^2,
    \label{eq:photon_relaxation_loaded_capacity}
\end{equation}
and the FCC nearest-neighbor separation is
\begin{equation}
    d_{\mathrm{nn}}
    =
    2r_{\mathrm{D}}.
    \label{eq:photon_relaxation_neighbor_spacing}
\end{equation}

The first objective is to construct a positive reduced lattice closure that preserves the continuum capacity \(c_{\mathrm{P}}\) while cancelling the leading quadratic group-speed correction of the nearest-neighbor branch. The second is to identify the additional constitutive conditions under which weak reservoir coupling produces coherent frequency transport rather than ordinary attenuation.

The logical status is:

\begin{enumerate}
    \item FCC geometry, primitive-cell inertia, coherent Condensate entrainment, and the long-wavelength endpoint-transfer capacity are inherited TFT structures.

    \item The dispersion-suppressed transverse branch is derived from the positive reduced action selected in this appendix.

    \item Uniform extension of the endpoint-transfer factor to finite wave number, the inverse-frequency loss tangent, modal-action preservation, and wave-vector transport are constitutive closures.

    \item Their recovery from the complete contact Hessian, rotational variables, internal phases, Condensate response, and Tension-Column dynamics remains a microscopic completion requirement.
\end{enumerate}

Accordingly, neither the coefficient \(1/12\) nor the nonconservative relaxation kernel is claimed to be uniquely derived from the complete FCC--Condensate parent theory. The conservative carrier branch is developed first; relaxation is introduced only after its finite-wave-number structure has been specified.

\section{Baseline Directional Branch}
\label{sec:photon_relaxation_baseline_branch}

Consider the reduced nearest-neighbor transfer chain associated with one FCC bond direction, with scalar wave number \(k\) measured along that direction. The gravity-neutral branch established in \hyperref[app:fcc_microphysics]{Appendix~\ref*{app:fcc_microphysics}} is
\begin{equation}
    \omega^2
    =
    \frac{4c_{\mathrm{E}}^2}{d_{\mathrm{nn}}^2}
    \sin^2
    \left(
        \frac{kd_{\mathrm{nn}}}{2}
    \right).
    \label{eq:photon_relaxation_baseline_neutral_branch}
\end{equation}

Applying the finite-wave-number loading inheritance stated above gives
\begin{equation}
    \omega^2
    =
    \frac{4c_{\mathrm{P}}^2}{d_{\mathrm{nn}}^2}
    \sin^2
    \left(
        \frac{kd_{\mathrm{nn}}}{2}
    \right).
    \label{eq:photon_relaxation_baseline_dispersion}
\end{equation}
This substitution extends the established long-wavelength endpoint-transfer capacity to the reduced finite-\(k\) branch; it is not yet derived from the fully loaded FCC dynamic matrix.

For the positive-frequency branch with
\(0\leq kd_{\mathrm{nn}}\leq\pi\),
\begin{equation}
    \omega
    =
    \frac{2c_{\mathrm{P}}}{d_{\mathrm{nn}}}
    \sin
    \left(
        \frac{kd_{\mathrm{nn}}}{2}
    \right),
    \label{eq:photon_relaxation_baseline_positive_branch}
\end{equation}
and the directional group velocity is
\begin{equation}
    v_{\mathrm{g}}
    =
    \frac{\mathrm{d}\omega}{\mathrm{d}k}
    =
    c_{\mathrm{P}}
    \cos
    \left(
        \frac{kd_{\mathrm{nn}}}{2}
    \right).
    \label{eq:photon_relaxation_baseline_group_velocity}
\end{equation}

For \(kd_{\mathrm{nn}}\ll1\),
\begin{equation}
    \omega^2
    =
    c_{\mathrm{P}}^2k^2
    \left[
        1
        -
        \frac{k^2d_{\mathrm{nn}}^2}{12}
        +
        \frac{k^4d_{\mathrm{nn}}^4}{360}
        +
        \mathcal{O}
        \left(
            k^6d_{\mathrm{nn}}^6
        \right)
    \right],
    \label{eq:photon_relaxation_baseline_dispersion_expansion}
\end{equation}
while
\begin{equation}
    v_{\mathrm{g}}
    =
    c_{\mathrm{P}}
    \left[
        1
        -
        \frac{k^2d_{\mathrm{nn}}^2}{8}
        +
        \frac{k^4d_{\mathrm{nn}}^4}{384}
        +
        \mathcal{O}
        \left(
            k^6d_{\mathrm{nn}}^6
        \right)
    \right].
    \label{eq:photon_relaxation_baseline_group_velocity_expansion}
\end{equation}

The continuum limit is
\begin{equation}
    \lim_{k\rightarrow0}
    \frac{\omega}{k}
    =
    \lim_{k\rightarrow0}
    v_{\mathrm{g}}
    =
    c_{\mathrm{P}}.
    \label{eq:photon_relaxation_baseline_continuum_limit}
\end{equation}
The nearest-neighbor chain therefore recovers the required local propagation capacity but retains a leading group-speed correction proportional to \(k^2d_{\mathrm{nn}}^2\).

Equation~\eqref{eq:photon_relaxation_baseline_dispersion} is exact only for the selected directional reduction. It is not the completed three-dimensional photon eigenbranch and does not contain the rotational, internal-phase, or auxiliary transfer structures that may modify the finite-wave-number spectrum.

\section{Dimensionally Closed Reduced Transverse Cell Action}
\label{sec:photon_relaxation_minimal_closure}

Let \(\widehat{\mathbf{n}}_a\), \(a=1,\ldots,6\), denote the six undirected FCC nearest-neighbor directions defined in \hyperref[app:fcc_microphysics]{Appendix~\ref*{app:fcc_microphysics}}. One representative is retained from each opposite-direction pair, so summation over \(\mathbf{R}\) and \(a\) counts every nearest-neighbor bond once. These directions satisfy
\begin{equation}
    \sum_{a=1}^{6}
    \widehat{n}_{a,i}
    \widehat{n}_{a,j}
    =
    2\delta_{ij}.
    \label{eq:photon_relaxation_fcc_second_moment}
\end{equation}

For the divergence-free scaffold displacement
\(\mathbf{u}(\mathbf{R},t)\), define the forward nearest-neighbor difference
\begin{equation}
    \Delta_a\mathbf{u}(\mathbf{R},t)
    =
    \mathbf{u}
    \left(
        \mathbf{R}
        +
        d_{\mathrm{nn}}
        \widehat{\mathbf{n}}_a,
        t
    \right)
    -
    \mathbf{u}(\mathbf{R},t),
    \label{eq:photon_relaxation_first_difference}
\end{equation}
and the centered second difference
\begin{align}
    \Delta_a^{(2)}
    \mathbf{u}(\mathbf{R},t)
    &=
    \mathbf{u}
    \left(
        \mathbf{R}
        +
        d_{\mathrm{nn}}
        \widehat{\mathbf{n}}_a,
        t
    \right)
    -
    2\mathbf{u}(\mathbf{R},t)
    \nonumber\\
    &\quad
    +
    \mathbf{u}
    \left(
        \mathbf{R}
        -
        d_{\mathrm{nn}}
        \widehat{\mathbf{n}}_a,
        t
    \right).
    \label{eq:photon_relaxation_second_difference}
\end{align}
The notation
\(\bigl(\Delta_a^{(2)}\bigr)^2\mathbf{u}\)
denotes repeated action of the centered operator, not the square of the forward difference.

The unstrained primitive-cell scaffold mass is
\begin{equation}
    \rho_{\mathrm{sh}}^{(\mathrm{micro},0)}
    V_{\mathrm{p}},
    \label{eq:photon_relaxation_cell_mass}
\end{equation}
where
\(\rho_{\mathrm{sh}}^{(\mathrm{micro},0)}\)
and \(V_{\mathrm{p}}\) are inherited from
\hyperref[app:fcc_microphysics]{Appendix~\ref*{app:fcc_microphysics}}.
Under the finite-wave-number endpoint-transfer inheritance, the referred translational inertia of the loaded transverse branch is
\begin{equation}
    \frac{
        \rho_{\mathrm{sh}}^{(\mathrm{micro},0)}
        V_{\mathrm{p}}
    }{
        \left(
            1-\alpha_{\mathrm{G}}
        \right)^4
    }.
    \label{eq:photon_relaxation_loaded_cell_inertia}
\end{equation}
This is an effective modal inertia, not a change in the underlying primitive-cell mass.

The reduced transverse lattice Lagrangian is
\begin{align}
    L_{\gamma}^{(\mathrm{T})}
    &=
    \sum_{\mathbf{R}}
    \Bigg\{
        \frac{1}{2}
        \frac{
            \rho_{\mathrm{sh}}^{(\mathrm{micro},0)}
            V_{\mathrm{p}}
        }{
            \left(
                1-\alpha_{\mathrm{G}}
            \right)^4
        }
        \left|
            \dot{\mathbf{u}}(\mathbf{R},t)
        \right|^2
    \nonumber\\
    &\qquad
        -
        \frac{
            \rho_{\mathrm{sh}}^{(\mathrm{micro},0)}
            V_{\mathrm{p}}
            c_{\mathrm{E}}^2
        }{
            4d_{\mathrm{nn}}^2
        }
        \sum_{a=1}^{6}
        \left[
            \left|
                \Delta_a\mathbf{u}
            \right|^2
            +
            \frac{1}{12}
            \left|
                \Delta_a^{(2)}\mathbf{u}
            \right|^2
        \right]
    \Bigg\}.
    \label{eq:photon_relaxation_cell_lagrangian}
\end{align}
The corresponding action is
\begin{equation}
    S_{\gamma}^{(\mathrm{T})}
    =
    \int
    L_{\gamma}^{(\mathrm{T})}
    \,\mathrm{d}t.
    \label{eq:photon_relaxation_reduced_action}
\end{equation}
Its variation is restricted to the coherent transverse subspace:
\begin{equation}
    \delta S_{\gamma}^{(\mathrm{T})}
    =
    0,
    \qquad
    \nabla\cdot\mathbf{u}
    =
    0,
    \qquad
    \nabla\cdot\delta\mathbf{u}
    =
    0.
    \label{eq:photon_relaxation_restricted_variation}
\end{equation}
Thus, Equation~\eqref{eq:photon_relaxation_cell_lagrangian} defines a reduced transverse carrier action rather than an unconstrained three-polarization lattice action.

Because both difference operators have dimensions of displacement,
\[
    \left[
        \rho_{\mathrm{sh}}^{(\mathrm{micro},0)}
        V_{\mathrm{p}}
        \left|
            \dot{\mathbf{u}}
        \right|^2
    \right]
    =
    \left[
        \rho_{\mathrm{sh}}^{(\mathrm{micro},0)}
        V_{\mathrm{p}}
        \frac{c_{\mathrm{E}}^2}{d_{\mathrm{nn}}^2}
        \left|
            \Delta_a\mathbf{u}
        \right|^2
    \right]
    =
    \left[
        \rho_{\mathrm{sh}}^{(\mathrm{micro},0)}
        V_{\mathrm{p}}
        \frac{c_{\mathrm{E}}^2}{d_{\mathrm{nn}}^2}
        \left|
            \Delta_a^{(2)}\mathbf{u}
        \right|^2
    \right]
    =
    \mathrm{J}.
\]
Hence
\([L_{\gamma}^{(\mathrm{T})}]=\mathrm{J}\)
and
\([S_{\gamma}^{(\mathrm{T})}]=\mathrm{J\,s}\).

The potential energy is nonnegative and vanishes under uniform translation. Together with
\(0\leq\alpha_{\mathrm{G}}<1\),
this ensures positive modal inertia and excludes a negative-stiffness sector within the selected transverse reduction.

Restricted variation of Equation~\eqref{eq:photon_relaxation_cell_lagrangian} gives
\begin{align}
    \ddot{\mathbf{u}}(\mathbf{R},t)
    &=
    \frac{
        c_{\mathrm{E}}^2
        \left(
            1-\alpha_{\mathrm{G}}
        \right)^4
    }{
        2d_{\mathrm{nn}}^2
    }
    \sum_{a=1}^{6}
    \left[
        \Delta_a^{(2)}\mathbf{u}
        -
        \frac{1}{12}
        \left(
            \Delta_a^{(2)}
        \right)^2
        \mathbf{u}
    \right]
    \nonumber\\
    &=
    \frac{
        c_{\mathrm{P}}^2
    }{
        2d_{\mathrm{nn}}^2
    }
    \sum_{a=1}^{6}
    \left[
        \Delta_a^{(2)}\mathbf{u}
        -
        \frac{1}{12}
        \left(
            \Delta_a^{(2)}
        \right)^2
        \mathbf{u}
    \right].
    \label{eq:photon_relaxation_extended_lattice_equation}
\end{align}

For the transverse Bloch mode
\begin{equation}
    \mathbf{u}(\mathbf{R},t)
    =
    \mathbf{u}_0
    \exp
    \left[
        i
        \left(
            \mathbf{k}\cdot\mathbf{R}
            -
            \omega t
        \right)
    \right],
    \qquad
    \mathbf{k}\cdot\mathbf{u}_0
    =
    0,
    \label{eq:photon_relaxation_extended_bloch_mode}
\end{equation}
the centered-difference symbol is
\begin{equation}
    \Delta_a^{(2)}
    \mathbf{u}
    =
    -4
    \sin^2
    \left(
        \frac{
            d_{\mathrm{nn}}
            \mathbf{k}\cdot
            \widehat{\mathbf{n}}_a
        }{
            2
        }
    \right)
    \mathbf{u}.
    \label{eq:photon_relaxation_second_difference_symbol}
\end{equation}

Writing
\[
    x
    =
    d_{\mathrm{nn}}
    \mathbf{k}\cdot
    \widehat{\mathbf{n}}_a,
\]
the forward- and centered-difference contributions are proportional to
\(4\sin^2(x/2)\)
and
\(16\sin^4(x/2)\), respectively. The normalized directional stiffness symbol is therefore
\begin{equation}
    \sin^2
    \left(
        \frac{x}{2}
    \right)
    +
    \frac{1}{3}
    \sin^4
    \left(
        \frac{x}{2}
    \right)
    =
    \frac{x^2}{4}
    -
    \frac{x^6}{360}
    +
    \frac{x^8}{4032}
    +
    \mathcal{O}
    \left(
        x^{10}
    \right).
    \label{eq:photon_relaxation_suppressed_symbol_expansion}
\end{equation}
The \(x^4\) term cancels exactly.

The real-space coefficient \(1/12\) is required because the squared centered difference contributes four times the normalized
\(\sin^4(x/2)\) symbol:
\[
    \frac{1}{12}\times4
    =
    \frac{1}{3}.
\]
This coefficient is unique within the selected positive two-term stencil. Longer-range contacts, rotational variables, internal-phase transfer, or eliminated auxiliary modes may reproduce the same low-wave-number cancellation through a different microscopic action.

\section{Exact Reduced Transverse Bloch Branch}
\label{sec:photon_relaxation_exact_branch}

Substitution of Equation~\eqref{eq:photon_relaxation_extended_bloch_mode} into Equation~\eqref{eq:photon_relaxation_extended_lattice_equation} gives
\begin{align}
    \omega^2
    \left(
        \mathbf{k},
        \alpha_{\mathrm{G}}
    \right)
    &=
    \frac{
        2c_{\mathrm{P}}^2
    }{
        d_{\mathrm{nn}}^2
    }
    \sum_{a=1}^{6}
    \Bigg[
        \sin^2
        \left(
            \frac{
                d_{\mathrm{nn}}
                \mathbf{k}\cdot
                \widehat{\mathbf{n}}_a
            }{
                2
            }
        \right)
    \nonumber\\
    &\hspace{7em}
        +
        \frac{1}{3}
        \sin^4
        \left(
            \frac{
                d_{\mathrm{nn}}
                \mathbf{k}\cdot
                \widehat{\mathbf{n}}_a
            }{
                2
            }
        \right)
    \Bigg].
    \label{eq:photon_relaxation_exact_bloch_branch}
\end{align}

Every term in Equation~\eqref{eq:photon_relaxation_exact_bloch_branch} is nonnegative. The selected transverse branch is therefore spectrally nonnegative, gapless at
\(\mathbf{k}=\mathbf{0}\), and invariant under
\(\mathbf{k}\rightarrow-\mathbf{k}\).
Because its scalar Bloch symbol acts identically on both displacement components orthogonal to \(\mathbf{k}\), the two transverse components remain degenerate within the selected reduction.

For \(\omega>0\), the group-velocity vector is
\begin{equation}
    \mathbf{v}_{\mathrm{g}}
    =
    \nabla_{\mathbf{k}}\omega.
    \label{eq:photon_relaxation_group_velocity_definition}
\end{equation}
Differentiating Equation~\eqref{eq:photon_relaxation_exact_bloch_branch} gives
\begin{align}
    \mathbf{v}_{\mathrm{g}}
    &=
    \frac{
        c_{\mathrm{P}}^2
    }{
        2\omega d_{\mathrm{nn}}
    }
    \sum_{a=1}^{6}
    \sin
    \left(
        d_{\mathrm{nn}}
        \mathbf{k}\cdot
        \widehat{\mathbf{n}}_a
    \right)
    \nonumber\\
    &\qquad\times
    \left[
        1
        +
        \frac{2}{3}
        \sin^2
        \left(
            \frac{
                d_{\mathrm{nn}}
                \mathbf{k}\cdot
                \widehat{\mathbf{n}}_a
            }{
                2
            }
        \right)
    \right]
    \widehat{\mathbf{n}}_a.
    \label{eq:photon_relaxation_exact_group_velocity_vector}
\end{align}
At finite wave number, FCC anisotropy permits
\(\mathbf{v}_{\mathrm{g}}\) to deviate from the direction of
\(\mathbf{k}\); this deviation vanishes in the continuum limit.

Define
\begin{equation}
    k
    =
    \left|
        \mathbf{k}
    \right|,
    \qquad
    \widehat{\mathbf{k}}
    =
    \frac{
        \mathbf{k}
    }{
        k
    }.
    \label{eq:photon_relaxation_wavevector_direction}
\end{equation}
Using Equation~\eqref{eq:photon_relaxation_suppressed_symbol_expansion}, the long-wavelength expansion is
\begin{align}
    \omega^2
    &=
    c_{\mathrm{P}}^2
    \Bigg[
        k^2
        -
        \frac{
            d_{\mathrm{nn}}^4
        }{
            180
        }
        \sum_{a=1}^{6}
        \left(
            \mathbf{k}\cdot
            \widehat{\mathbf{n}}_a
        \right)^6
    \nonumber\\
    &\hspace{6em}
        +
        \frac{
            d_{\mathrm{nn}}^6
        }{
            2016
        }
        \sum_{a=1}^{6}
        \left(
            \mathbf{k}\cdot
            \widehat{\mathbf{n}}_a
        \right)^8
        +
        \mathcal{O}
        \left(
            k^{10}d_{\mathrm{nn}}^8
        \right)
    \Bigg].
    \label{eq:photon_relaxation_exact_branch_expansion}
\end{align}
The leading correction to \(\omega^2\) therefore begins at order
\(k^6d_{\mathrm{nn}}^4\), rather than
\(k^4d_{\mathrm{nn}}^2\).

Taking the positive square root gives
\begin{equation}
    \omega
    =
    c_{\mathrm{P}}k
    \left[
        1
        -
        \frac{
            k^4d_{\mathrm{nn}}^4
        }{
            360
        }
        \sum_{a=1}^{6}
        \left(
            \widehat{\mathbf{k}}
            \cdot
            \widehat{\mathbf{n}}_a
        \right)^6
        +
        \mathcal{O}
        \left(
            k^6d_{\mathrm{nn}}^6
        \right)
    \right].
    \label{eq:photon_relaxation_exact_frequency_expansion}
\end{equation}

For a fixed propagation direction, define the radial group speed by
\begin{equation}
    v_{\mathrm{g},\parallel}
    =
    \left.
        \frac{
            \partial\omega
        }{
            \partial k
        }
    \right|_{\widehat{\mathbf{k}}}
    =
    \mathbf{v}_{\mathrm{g}}
    \cdot
    \widehat{\mathbf{k}}.
    \label{eq:photon_relaxation_radial_group_speed}
\end{equation}
Its expansion is
\begin{equation}
    v_{\mathrm{g},\parallel}
    =
    c_{\mathrm{P}}
    \left[
        1
        -
        \frac{
            k^4d_{\mathrm{nn}}^4
        }{
            72
        }
        \sum_{a=1}^{6}
        \left(
            \widehat{\mathbf{k}}
            \cdot
            \widehat{\mathbf{n}}_a
        \right)^6
        +
        \mathcal{O}
        \left(
            k^6d_{\mathrm{nn}}^6
        \right)
    \right].
    \label{eq:photon_relaxation_suppressed_group_velocity}
\end{equation}

For the six undirected FCC nearest-neighbor directions, the cubic directional factor satisfies
\begin{equation}
    \frac{1}{2}
    \leq
    \sum_{a=1}^{6}
    \left(
        \widehat{\mathbf{k}}
        \cdot
        \widehat{\mathbf{n}}_a
    \right)^6
    \leq
    \frac{17}{16}.
    \label{eq:photon_relaxation_fcc_directional_bounds}
\end{equation}
Its principal-direction values are
\[
    \frac{1}{2}
    \quad\text{along}\quad
    [100],
    \qquad
    \frac{17}{16}
    \quad\text{along}\quad
    [110],
    \qquad
    \frac{8}{9}
    \quad\text{along}\quad
    [111].
\]

The leading correction to the radial group speed is therefore proportional to
\(k^4d_{\mathrm{nn}}^4\), replacing the
\(k^2d_{\mathrm{nn}}^2\) correction of the baseline chain. The continuum limits remain
\begin{equation}
    \lim_{k\rightarrow0}
    \frac{
        \omega
    }{
        k
    }
    =
    \lim_{k\rightarrow0}
    v_{\mathrm{g},\parallel}
    =
    c_{\mathrm{P}}.
    \label{eq:photon_relaxation_suppressed_continuum_limit}
\end{equation}

The selected action thus suppresses the leading nearest-neighbor dispersion without changing the endpoint-transfer capacity, opening a spectral gap, or introducing negative stiffness. These conclusions apply to the reduced transverse branch and do not constitute a stability proof for the complete multi-branch FCC spectrum.

\section{Conservative Condensate Limit and Separate Nonconservative Response}
\label{sec:photon_relaxation_condensate_completion}

The dispersion-suppressed branch derived in Section~\ref{sec:photon_relaxation_exact_branch} belongs to the conservative coherent sector of the dual-phase vacuum. The constraint inherited from \hyperref[app:condensate_impedance]{Appendix~\ref*{app:condensate_impedance}} requires
\begin{equation}
    \mathbf{v}_{\mathrm{rel}}
    =
    \mathbf{v}_{\mathrm{IC}}
    -
    \dot{\mathbf{u}}
    =
    \mathbf{0}
    \label{eq:photon_relaxation_zero_relative_velocity}
\end{equation}
at the retained carrier order. The FCC scaffold and interstitial Condensate therefore share the same incremental velocity, while the Condensate supplies the pressure response enforcing the transverse constraint.

Equation~\eqref{eq:photon_relaxation_zero_relative_velocity} excludes relative flow and dissipation from the leading conservative action but does not prohibit a weaker irreversible response at higher order. Vacuum relaxation must therefore enter through a distinct nonconservative sector rather than through additional carrier inertia or a multiplicative correction to \(c_{\mathrm{P}}\).

Let \(\widehat{\mathbf{e}}_{\mathrm{T}}\) be a normalized transverse polarization:
\begin{equation}
    \widehat{\mathbf{e}}_{\mathrm{T}}
    \cdot
    \mathbf{k}
    =
    0,
    \qquad
    \left|
        \widehat{\mathbf{e}}_{\mathrm{T}}
    \right|
    =
    1.
    \label{eq:photon_relaxation_transverse_polarization}
\end{equation}
Define the work-conjugate projected strain and stress by
\begin{equation}
    \varepsilon_{\mathrm{T}}
    =
    \widehat{e}_{\mathrm{T},i}
    \widehat{e}_{\mathrm{T},j}
    \varepsilon_{ij},
    \qquad
    \sigma_{\mathrm{T}}
    =
    \widehat{e}_{\mathrm{T},i}
    \widehat{e}_{\mathrm{T},j}
    \sigma_{ij}.
    \label{eq:photon_relaxation_modal_projection}
\end{equation}

For a locally harmonic component, adopt
\begin{equation}
    \varepsilon_{\mathrm{T}}(\mathbf{x},t)
    =
    \varepsilon_{\mathrm{T}}^{(0)}(\mathbf{x})
    \cos
    \left(
        \omega t
    \right),
    \label{eq:photon_relaxation_modal_strain}
\end{equation}
and
\begin{equation}
    \sigma_{\mathrm{T}}(\mathbf{x},t)
    =
    \sigma_{\mathrm{T}}^{(0)}(\mathbf{x})
    \cos
    \left(
        \omega t
        +
        \Delta\varphi
    \right).
    \label{eq:photon_relaxation_modal_stress}
\end{equation}
Under this convention,
\begin{equation}
    \Delta\varphi
    >
    0
    \label{eq:photon_relaxation_positive_phase_lag}
\end{equation}
means that the projected stress leads the strain and transfers positive work from the coherent mode to the reservoir. The conservative limit is
\begin{equation}
    \Delta\varphi
    =
    0,
    \label{eq:photon_relaxation_zero_phase_lag}
\end{equation}
for which the stress--strain loop encloses no area and the net irreversible work per cycle vanishes.

The retained response is weak and adiabatic:
\begin{equation}
    \frac{
        \Gamma_{\mathrm{D}}
    }{
        \omega
    }
    \ll
    1.
    \label{eq:photon_relaxation_weak_loss_condition}
\end{equation}
A nonzero phase lag can therefore establish a slow modal-energy transfer rate, but it does not by itself determine whether the response produces attenuation, absorption, mode conversion, or coherent frequency evolution.

\section{Modal Phase Lag and the Viscoelastic Vacuum Relaxation Rate}
\label{sec:photon_relaxation_phase_lag_rate}

Let \(\mathcal{V}_{\gamma}\) denote the instantaneous spatial support of the coherent mode. Over one optical cycle, the modal amplitudes and support may be treated as fixed. The positive energy transferred out of the carrier is
\begin{equation}
    \Delta E_{\mathrm{cycle}}
    =
    \int_{\mathcal{V}_{\gamma}}
    \oint_{\mathrm{cycle}}
    \sigma_{\mathrm{T}}\,
    \mathrm{d}\varepsilon_{\mathrm{T}}\,
    \mathrm{d}V.
    \label{eq:photon_relaxation_cycle_work_definition}
\end{equation}
Because strain is dimensionless,
\[
    \left[
        \sigma_{\mathrm{T}}\,
        \mathrm{d}\varepsilon_{\mathrm{T}}\,
        \mathrm{d}V
    \right]
    =
    \mathrm{Pa\,m^3}
    =
    \mathrm{J}.
\]

Using Equations~\eqref{eq:photon_relaxation_modal_strain} and~\eqref{eq:photon_relaxation_modal_stress},
\begin{equation}
    \Delta E_{\mathrm{cycle}}
    =
    \pi
    \int_{\mathcal{V}_{\gamma}}
    \sigma_{\mathrm{T}}^{(0)}
    \varepsilon_{\mathrm{T}}^{(0)}
    \sin
    \left(
        \Delta\varphi
    \right)
    \mathrm{d}V.
    \label{eq:photon_relaxation_cycle_work}
\end{equation}
The maximum recoverable energy associated with the in-phase response is
\begin{equation}
    E_{\mathrm{st}}
    =
    \frac{1}{2}
    \int_{\mathcal{V}_{\gamma}}
    \sigma_{\mathrm{T}}^{(0)}
    \varepsilon_{\mathrm{T}}^{(0)}
    \cos
    \left(
        \Delta\varphi
    \right)
    \mathrm{d}V.
    \label{eq:photon_relaxation_stored_modal_energy}
\end{equation}
For a spatially uniform phase lag,
\begin{equation}
    \frac{
        \Delta E_{\mathrm{cycle}}
    }{
        E_{\mathrm{st}}
    }
    =
    2\pi
    \tan
    \left(
        \Delta\varphi
    \right).
    \label{eq:photon_relaxation_fractional_cycle_loss}
\end{equation}

For a nonuniform response, define the effective modal loss tangent by
\begin{equation}
    \tan
    \left(
        \Delta\varphi_{\mathrm{eff}}
    \right)
    \equiv
    \frac{
        \displaystyle
        \int_{\mathcal{V}_{\gamma}}
        \sigma_{\mathrm{T}}^{(0)}
        \varepsilon_{\mathrm{T}}^{(0)}
        \sin
        \left(
            \Delta\varphi
        \right)
        \mathrm{d}V
    }{
        \displaystyle
        \int_{\mathcal{V}_{\gamma}}
        \sigma_{\mathrm{T}}^{(0)}
        \varepsilon_{\mathrm{T}}^{(0)}
        \cos
        \left(
            \Delta\varphi
        \right)
        \mathrm{d}V
    }.
    \label{eq:photon_relaxation_effective_loss_tangent}
\end{equation}
Equation~\eqref{eq:photon_relaxation_fractional_cycle_loss} then remains valid with
\(\Delta\varphi_{\mathrm{eff}}\), denoted simply by
\(\Delta\varphi\) below.

Let the slowly varying coherent modal energy obey
\begin{equation}
    \frac{
        \mathrm{d}E
    }{
        \mathrm{d}t
    }
    =
    -\Gamma_{\mathrm{D}}E.
    \label{eq:photon_relaxation_energy_decay_reintroduced}
\end{equation}
The coefficient $\Gamma_{\mathrm{D}}$ is termed the \emph{Viscoelastic Vacuum Relaxation Rate}, or equivalently the \emph{FCC--Condensate relaxation rate}. It is the macroscopic fractional coherent-energy relaxation rate associated with the weak out-of-phase viscoelastic response of the FCC--Condensate vacuum to the propagating transverse photon mode. The corresponding stress--strain phase lag provides the nonconservative work channel through which energy is transferred from the coherent carrier into non-coherent FCC--Condensate reservoir degrees of freedom.
In the weak-loss regime,
\begin{equation}
    E
    =
    E_{\mathrm{st}}
    \left[
        1
        +
        \mathcal{O}
        \left(
            \frac{
                \Gamma_{\mathrm{D}}^2
            }{
                \omega^2
            }
        \right)
    \right],
    \label{eq:photon_relaxation_total_stored_energy_relation}
\end{equation}
so replacing \(E_{\mathrm{st}}\) by \(E\) does not change the leading-order result.

Over one period \(2\pi/\omega\), Equation~\eqref{eq:photon_relaxation_energy_decay_reintroduced} gives
\begin{equation}
    \frac{
        \Delta E_{\mathrm{cycle}}
    }{
        E
    }
    =
    1
    -
    \exp
    \left(
        -
        \frac{
            2\pi\Gamma_{\mathrm{D}}
        }{
            \omega
        }
    \right).
    \label{eq:photon_relaxation_exact_cycle_decay}
\end{equation}
Expanding under Equation~\eqref{eq:photon_relaxation_weak_loss_condition},
\begin{equation}
    \frac{
        \Delta E_{\mathrm{cycle}}
    }{
        E
    }
    =
    \frac{
        2\pi\Gamma_{\mathrm{D}}
    }{
        \omega
    }
    +
    \mathcal{O}
    \left(
        \frac{
            \Gamma_{\mathrm{D}}^2
        }{
            \omega^2
        }
    \right).
    \label{eq:photon_relaxation_cycle_decay_expansion}
\end{equation}

Equating the stress--strain work with the temporal energy loss yields
\begin{equation}
    \Gamma_{\mathrm{D}}
    =
    \omega
    \tan
    \left(
        \Delta\varphi
    \right)
    +
    \mathcal{O}
    \left(
        \frac{
            \Gamma_{\mathrm{D}}^2
        }{
            \omega
        }
    \right),
    \label{eq:photon_relaxation_gamma_phase_relation}
\end{equation}
and therefore
\begin{equation}
    \Gamma_{\mathrm{D}}
    \simeq
    \omega
    \tan
    \left(
        \Delta\varphi
    \right).
    \label{eq:photon_relaxation_gamma_phase_leading}
\end{equation}

An approximately frequency-independent fractional transfer rate requires
\begin{equation}
    \tan
    \left[
        \Delta\varphi(\omega)
    \right]
    \simeq
    \frac{
        \Gamma_{\mathrm{D}}
    }{
        \omega
    },
    \label{eq:photon_relaxation_inverse_frequency_phase_lag}
\end{equation}
or, for a small lag,
\begin{equation}
    \Delta\varphi(\omega)
    \simeq
    \frac{
        \Gamma_{\mathrm{D}}
    }{
        \omega
    }.
    \label{eq:photon_relaxation_small_phase_lag}
\end{equation}
The corresponding fractional transfer per cycle is
\begin{equation}
    \frac{
        \Delta E_{\mathrm{cycle}}
    }{
        E
    }
    \simeq
    \frac{
        2\pi\Gamma_{\mathrm{D}}
    }{
        \omega
    }.
    \label{eq:photon_relaxation_cycle_loss_inverse_frequency}
\end{equation}
Higher-frequency modes lose a smaller fraction during each cycle but complete proportionally more cycles per unit time.

A frequency-independent phase lag would instead imply
\begin{equation}
    \Gamma_{\mathrm{D}}
    \propto
    \omega,
    \label{eq:photon_relaxation_constant_phase_chromatic}
\end{equation}
and hence a chromatic transfer rate. The inverse-frequency loss tangent is therefore necessary for the proposed achromatic relaxation law.

This section determines the response required by a prescribed constant
\(\Gamma_{\mathrm{D}}\); it does not derive that response from FCC--Condensate microphysics. The parent theory must provide a causal susceptibility whose reactive and dissipative components satisfy the corresponding Kramers--Kronig relations and whose reservoir coupling remains consistent with phase noise, line broadening, and fluctuation--dissipation constraints.

Most importantly, the phase lag establishes an energy-transfer rate only. In a stationary linear medium, such a response ordinarily produces attenuation or a complex wave number, not secular motion through the real-frequency eigenbranch.

\section{Ordinary Linear Damping and the Amplitude-Attenuation No-Go}
\label{sec:photon_relaxation_attenuation_no_go}

To isolate this distinction, consider one fixed transverse Fourier mode with normal coordinate \(q_{\mathbf{k}}(t)\) and conservative eigenfrequency
\(\omega_{\mathrm{c}}(\mathbf{k})\) from Equation~\eqref{eq:photon_relaxation_exact_bloch_branch}. Conventional linear damping gives
\begin{equation}
    \ddot{q}_{\mathbf{k}}
    +
    \Gamma_{\mathrm{D}}
    \dot{q}_{\mathbf{k}}
    +
    \omega_{\mathrm{c}}^2(\mathbf{k})
    q_{\mathbf{k}}
    =
    0.
    \label{eq:photon_relaxation_linearly_damped_mode}
\end{equation}
This equation is a diagnostic of ordinary damping, not the completed TFT transport law.

For
\begin{equation}
    \omega_{\mathrm{c}}^2(\mathbf{k})
    >
    \frac{
        \Gamma_{\mathrm{D}}^2
    }{
        4
    },
    \label{eq:photon_relaxation_underdamped_condition}
\end{equation}
the solution is
\begin{equation}
    q_{\mathbf{k}}(t)
    =
    q_{\mathbf{k}}^{(0)}
    \exp
    \left(
        -
        \frac{
            \Gamma_{\mathrm{D}}t
        }{
            2
        }
    \right)
    \exp
    \left(
        -i\Omega_{\mathbf{k}}t
    \right),
    \label{eq:photon_relaxation_damped_mode_solution}
\end{equation}
where
\begin{equation}
    \Omega_{\mathbf{k}}
    =
    \sqrt{
        \omega_{\mathrm{c}}^2(\mathbf{k})
        -
        \frac{
            \Gamma_{\mathrm{D}}^2
        }{
            4
        }
    }.
    \label{eq:photon_relaxation_damped_frequency_definition}
\end{equation}

The amplitude and cycle-averaged energy therefore scale as
\begin{equation}
    \left|
        q_{\mathbf{k}}(t)
    \right|
    \propto
    \exp
    \left(
        -
        \frac{
            \Gamma_{\mathrm{D}}t
        }{
            2
        }
    \right),
    \label{eq:photon_relaxation_amplitude_attenuation}
\end{equation}
and
\begin{equation}
    E_{\mathbf{k}}(t)
    \propto
    \exp
    \left(
        -
        \Gamma_{\mathrm{D}}t
    \right).
    \label{eq:photon_relaxation_linear_attenuation}
\end{equation}

The real oscillation frequency is
\begin{align}
    \Omega_{\mathbf{k}}
    &=
    \omega_{\mathrm{c}}(\mathbf{k})
    -
    \frac{
        \Gamma_{\mathrm{D}}^2
    }{
        8\omega_{\mathrm{c}}(\mathbf{k})
    }
    +
    \mathcal{O}
    \left(
        \frac{
            \Gamma_{\mathrm{D}}^4
        }{
            \omega_{\mathrm{c}}^3(\mathbf{k})
        }
    \right).
    \label{eq:photon_relaxation_damped_real_frequency}
\end{align}
For fixed \(\mathbf{k}\), this is a constant second-order shift, not a secular law of the form
\begin{equation}
    \frac{
        \mathrm{d}\omega
    }{
        \mathrm{d}t
    }
    =
    -\Gamma_{\mathrm{D}}\omega.
    \label{eq:photon_relaxation_frequency_drift_target}
\end{equation}

The equivalent spatial description uses a complex wave number,
\begin{equation}
    k
    =
    k_{\mathrm{R}}
    +
    ik_{\mathrm{I}},
    \qquad
    k_{\mathrm{I}}
    >
    0,
    \label{eq:photon_relaxation_complex_wave_number}
\end{equation}
giving
\begin{equation}
    q(\ell,t)
    \propto
    \exp
    \left(
        -
        k_{\mathrm{I}}\ell
    \right)
    \exp
    \left[
        i
        \left(
            k_{\mathrm{R}}\ell
            -
            \omega t
        \right)
    \right].
    \label{eq:photon_relaxation_spatial_attenuation}
\end{equation}
The amplitude decreases while the imposed real frequency remains fixed.

Ordinary linear damping can therefore describe amplitude decay, intensity attenuation, absorption, or loss of coherent-mode survival probability. None of these outcomes alone establishes a continuous redshift of the surviving coherent carrier. Consequently,
\begin{equation}
    \frac{
        \mathrm{d}E
    }{
        \mathrm{d}t
    }
    =
    -\Gamma_{\mathrm{D}}E
    \quad
    \not\Rightarrow
    \quad
    \frac{
        \mathrm{d}\omega
    }{
        \mathrm{d}t
    }
    =
    -\Gamma_{\mathrm{D}}\omega
    \label{eq:photon_relaxation_attenuation_no_go}
\end{equation}
without an additional condition that moves the coherent mode through the real-frequency eigenstate family.

The required completion must therefore be a coherent inelastic transport channel distinct from ordinary linear attenuation. It must change the real carrier frequency and wave vector while suppressing incoherent absorption, transverse diffusion, polarization scrambling, phase noise, and secondary-mode radiation.

\section{Coherent-Action-Preserving Frequency Transport}
\label{sec:photon_relaxation_frequency_transport}

Section~\ref{sec:photon_relaxation_attenuation_no_go} shows that the modal-energy law
\begin{equation}
    \frac{\mathrm{d}E}{\mathrm{d}t}
    =
    -\Gamma_{\mathrm{D}}E
    \label{eq:photon_relaxation_transport_energy_law}
\end{equation}
does not determine the evolution of the real carrier frequency. The additional constitutive condition adopted here is preservation of the coherent modal action
\begin{equation}
    \mathcal{I}_{\gamma}
    \equiv
    \frac{E}{\omega},
    \label{eq:photon_relaxation_coherent_modal_action}
\end{equation}
with
\begin{equation}
    \left[
        \mathcal{I}_{\gamma}
    \right]
    =
    \mathrm{J\,s}.
    \label{eq:photon_relaxation_coherent_modal_action_dimensions}
\end{equation}

During the weak coherent-relaxation interval,
\begin{equation}
    \frac{\mathrm{d}\mathcal{I}_{\gamma}}{\mathrm{d}t}
    =
    \frac{\mathrm{d}}{\mathrm{d}t}
    \left(
        \frac{E}{\omega}
    \right)
    =
    0.
    \label{eq:photon_relaxation_wave_action_invariance}
\end{equation}
This condition applies to a coherent finite packet and includes a quantum-normalized excitation as the special case
\(\mathcal{I}_{\gamma}=\hbar\). It requires preservation of the selected modal identity and excludes incoherent conversion into other branches at the retained order.

Expanding Equation~\eqref{eq:photon_relaxation_wave_action_invariance} gives
\begin{equation}
    \frac{1}{\omega}
    \frac{\mathrm{d}E}{\mathrm{d}t}
    -
    \frac{E}{\omega^2}
    \frac{\mathrm{d}\omega}{\mathrm{d}t}
    =
    0,
    \label{eq:photon_relaxation_wave_action_expansion}
\end{equation}
or
\begin{equation}
    \frac{1}{E}
    \frac{\mathrm{d}E}{\mathrm{d}t}
    =
    \frac{1}{\omega}
    \frac{\mathrm{d}\omega}{\mathrm{d}t}.
    \label{eq:photon_relaxation_equal_fractional_rates}
\end{equation}
Combining this relation with Equation~\eqref{eq:photon_relaxation_transport_energy_law} yields
\begin{equation}
    \frac{\mathrm{d}\omega}{\mathrm{d}t}
    =
    -\Gamma_{\mathrm{D}}\omega.
    \label{eq:photon_relaxation_frequency_transport_law}
\end{equation}

The conditional transport logic is
\begin{equation}
    \begin{aligned}
        &\text{modal-energy transfer}
        +
        \text{coherent-action preservation}
        \\
        &\hspace{6em}
        \longrightarrow
        \text{real-frequency transport}.
    \end{aligned}
    \label{eq:photon_relaxation_frequency_transport_logic}
\end{equation}
Thus, Equation~\eqref{eq:photon_relaxation_frequency_transport_law} follows from the combined closure, not from ordinary linear damping alone.

For constant \(\Gamma_{\mathrm{D}}\),
\begin{equation}
    \omega(t)
    =
    \omega_{\mathrm{emit}}
    \exp
    \left[
        -
        \Gamma_{\mathrm{D}}
        \left(
            t-t_{\mathrm{emit}}
        \right)
    \right],
    \label{eq:photon_relaxation_frequency_solution}
\end{equation}
and
\begin{equation}
    E(t)
    =
    E_{\mathrm{emit}}
    \exp
    \left[
        -
        \Gamma_{\mathrm{D}}
        \left(
            t-t_{\mathrm{emit}}
        \right)
    \right].
    \label{eq:photon_relaxation_transport_energy_solution}
\end{equation}
Their ratio remains invariant:
\begin{equation}
    \frac{E(t)}{\omega(t)}
    =
    \frac{E_{\mathrm{emit}}}{\omega_{\mathrm{emit}}}
    =
    \mathcal{I}_{\gamma}.
    \label{eq:photon_relaxation_action_ratio_preserved}
\end{equation}

\subsection{Required Wave-Vector-Space Transport}
\label{subsec:photon_relaxation_wavevector_drag}

Frequency evolution requires motion through the real-frequency eigenbranch. Let
\begin{equation}
    \omega
    =
    \omega_{\mathrm{c}}
    \left(
        \mathbf{x},
        \mathbf{k}
    \right)
    \label{eq:photon_relaxation_local_conservative_branch}
\end{equation}
denote the local conservative branch, including quasistatic spatial dependence through
\(\alpha_{\mathrm{G}}(\mathbf{x})\). For a stationary background,
\begin{equation}
    \frac{\partial\omega_{\mathrm{c}}}{\partial t}
    =
    0.
    \label{eq:photon_relaxation_stationary_branch}
\end{equation}

Extend the conservative ray equations by a nonconservative wave-vector term:
\begin{equation}
    \dot{\mathbf{x}}
    =
    \nabla_{\mathbf{k}}\omega_{\mathrm{c}}
    \equiv
    \mathbf{v}_{\mathrm{g}},
    \label{eq:photon_relaxation_ray_position_equation}
\end{equation}
\begin{equation}
    \dot{\mathbf{k}}
    =
    -
    \nabla_{\mathbf{x}}\omega_{\mathrm{c}}
    +
    \left.
        \dot{\mathbf{k}}
    \right|_{\mathrm{D}}.
    \label{eq:photon_relaxation_ray_wavevector_equation}
\end{equation}
The first term in Equation~\eqref{eq:photon_relaxation_ray_wavevector_equation} produces conservative refraction; the second produces irreversible transport through wave-vector space.

Along the ray,
\begin{align}
    \frac{\mathrm{d}\omega_{\mathrm{c}}}{\mathrm{d}t}
    &=
    \nabla_{\mathbf{x}}\omega_{\mathrm{c}}
    \cdot
    \dot{\mathbf{x}}
    +
    \nabla_{\mathbf{k}}\omega_{\mathrm{c}}
    \cdot
    \dot{\mathbf{k}}
    \nonumber\\
    &=
    \mathbf{v}_{\mathrm{g}}
    \cdot
    \left.
        \dot{\mathbf{k}}
    \right|_{\mathrm{D}},
    \label{eq:photon_relaxation_frequency_ray_derivative}
\end{align}
because the conservative Hamiltonian contributions cancel.

Define the radial group-speed component on the outward branch by
\begin{equation}
    v_{\mathrm{g},\parallel}
    =
    \mathbf{v}_{\mathrm{g}}
    \cdot
    \widehat{\mathbf{k}}
    >
    0.
    \label{eq:photon_relaxation_positive_radial_group_speed}
\end{equation}
The minimal direction-preserving closure reproducing Equation~\eqref{eq:photon_relaxation_frequency_transport_law} is
\begin{equation}
    \left.
        \dot{\mathbf{k}}
    \right|_{\mathrm{D}}
    =
    -
    \Gamma_{\mathrm{D}}
    \frac{
        \omega_{\mathrm{c}}
    }{
        v_{\mathrm{g},\parallel}
    }
    \widehat{\mathbf{k}}.
    \label{eq:photon_relaxation_longitudinal_k_drag}
\end{equation}
Its dimensions are
\begin{equation}
    \left[
        \left.
            \dot{\mathbf{k}}
        \right|_{\mathrm{D}}
    \right]
    =
    \mathrm{m^{-1}\,s^{-1}}.
    \label{eq:photon_relaxation_wavevector_drag_dimensions}
\end{equation}
Substitution into Equation~\eqref{eq:photon_relaxation_frequency_ray_derivative} gives
\begin{equation}
    \frac{\mathrm{d}\omega_{\mathrm{c}}}{\mathrm{d}t}
    =
    -\Gamma_{\mathrm{D}}
    \omega_{\mathrm{c}}.
    \label{eq:photon_relaxation_k_drag_frequency_result}
\end{equation}

Because
\begin{equation}
    \widehat{\mathbf{k}}
    \times
    \left.
        \dot{\mathbf{k}}
    \right|_{\mathrm{D}}
    =
    \mathbf{0},
    \label{eq:photon_relaxation_no_dissipative_deflection}
\end{equation}
the dissipative term changes the wave-number magnitude without producing deterministic angular deflection. Conservative refraction remains governed by
\(-\nabla_{\mathbf{x}}\omega_{\mathrm{c}}\), while suppression of stochastic angular diffusion remains a microscopic requirement.

For propagation along a fixed direction,
\begin{equation}
    \left.
        \frac{\mathrm{d}k}{\mathrm{d}t}
    \right|_{\mathrm{D}}
    =
    -
    \Gamma_{\mathrm{D}}
    \frac{
        \omega_{\mathrm{c}}(k)
    }{
        v_{\mathrm{g},\parallel}(k)
    }.
    \label{eq:photon_relaxation_scalar_wavevector_drag}
\end{equation}
In the long-wavelength regime,
\begin{equation}
    \omega_{\mathrm{c}}
    \simeq
    c_{\mathrm{P}}k,
    \qquad
    v_{\mathrm{g},\parallel}
    \simeq
    c_{\mathrm{P}},
    \label{eq:photon_relaxation_long_wave_transport_limit}
\end{equation}
and therefore
\begin{equation}
    \frac{\mathrm{d}k}{\mathrm{d}t}
    \simeq
    -\Gamma_{\mathrm{D}}k.
    \label{eq:photon_relaxation_wave_number_decay}
\end{equation}
For constant \(\Gamma_{\mathrm{D}}\),
\begin{equation}
    k(t)
    =
    k_{\mathrm{emit}}
    \exp
    \left[
        -
        \Gamma_{\mathrm{D}}
        \left(
            t-t_{\mathrm{emit}}
        \right)
    \right],
    \label{eq:photon_relaxation_wave_number_solution}
\end{equation}
and, since \(\lambda=2\pi/k\),
\begin{equation}
    \lambda(t)
    =
    \lambda_{\mathrm{emit}}
    \exp
    \left[
        \Gamma_{\mathrm{D}}
        \left(
            t-t_{\mathrm{emit}}
        \right)
    \right].
    \label{eq:photon_relaxation_wavelength_solution}
\end{equation}

The wavelength increase is therefore the reciprocal wave-number response required for the carrier to remain on the real-frequency FCC eigenbranch. Equation~\eqref{eq:photon_relaxation_longitudinal_k_drag} remains a constitutive closure until it is recovered from the same causal FCC--Condensate response that produces the modal phase lag.

\subsection{Canonical Wave-Momentum and Reservoir Balance}
\label{subsec:photon_relaxation_reservoir_balance}

The coherent packet carries the canonical Bloch momentum
\begin{equation}
    \mathbf{p}_{\gamma}^{(\mathrm{can})}
    =
    \mathcal{I}_{\gamma}\mathbf{k}
    =
    \frac{E}{\omega}\mathbf{k}.
    \label{eq:photon_relaxation_packet_momentum}
\end{equation}
Its dimensions are
\begin{equation}
    \left[
        \mathbf{p}_{\gamma}^{(\mathrm{can})}
    \right]
    =
    \mathrm{kg\,m\,s^{-1}}.
    \label{eq:photon_relaxation_packet_momentum_dimensions}
\end{equation}

Since
\(\mathrm{d}\mathcal{I}_{\gamma}/\mathrm{d}t=0\),
the dissipative canonical-momentum change is
\begin{equation}
    \left.
        \frac{\mathrm{d}
        \mathbf{p}_{\gamma}^{(\mathrm{can})}}
        {\mathrm{d}t}
    \right|_{\mathrm{D}}
    =
    \mathcal{I}_{\gamma}
    \left.
        \dot{\mathbf{k}}
    \right|_{\mathrm{D}}.
    \label{eq:photon_relaxation_packet_momentum_loss}
\end{equation}

In the absence of additional loss channels, the corresponding reservoir bookkeeping is
\begin{equation}
    \frac{\mathrm{d}E_{\mathrm{res}}}{\mathrm{d}t}
    =
    \Gamma_{\mathrm{D}}E,
    \qquad
    \left.
        \frac{\mathrm{d}
        \mathbf{p}_{\mathrm{res}}^{(\mathrm{can})}}
        {\mathrm{d}t}
    \right|_{\mathrm{D}}
    =
    -
    \left.
        \frac{\mathrm{d}
        \mathbf{p}_{\gamma}^{(\mathrm{can})}}
        {\mathrm{d}t}
    \right|_{\mathrm{D}}.
    \label{eq:photon_relaxation_reservoir_balance}
\end{equation}

This is an energy--quasimomentum balance, not yet a derivation of mechanical substrate recoil. Because \(\mathbf{k}\) is a Bloch wave vector, the mapping
\begin{equation}
    \text{canonical lattice quasimomentum}
    \longrightarrow
    \text{mechanical substrate momentum}
    \label{eq:photon_relaxation_quasimomentum_mechanical_mapping}
\end{equation}
must be recovered from the FCC--Condensate microphysics. The same completion must constrain recoil radiation, incoherent mode conversion, image diffusion, phase noise, and polarization scrambling.

\section{Exact Relaxation Redshift and Euclidean Path Law}
\label{sec:photon_relaxation_exact_redshift}

Equation~\eqref{eq:photon_relaxation_frequency_transport_law} may be written as
\begin{equation}
    \frac{\mathrm{d}}{\mathrm{d}t}
    \ln\omega
    =
    -\Gamma_{\mathrm{D}}.
    \label{eq:photon_relaxation_log_frequency_transport}
\end{equation}
For a generally position- or time-dependent relaxation rate,
\begin{equation}
    \frac{
        \omega_{\mathrm{obs}}
    }{
        \omega_{\mathrm{emit}}
    }
    =
    \exp
    \left[
        -
        \int_{t_{\mathrm{emit}}}^{t_{\mathrm{obs}}}
        \Gamma_{\mathrm{D}}(t)
        \,\mathrm{d}t
    \right].
    \label{eq:photon_relaxation_frequency_ratio}
\end{equation}
Using
\begin{equation}
    1+z
    =
    \frac{
        \omega_{\mathrm{emit}}
    }{
        \omega_{\mathrm{obs}}
    },
    \label{eq:photon_relaxation_redshift_definition}
\end{equation}
gives the general temporal relation
\begin{equation}
    \ln
    \left(
        1+z
    \right)
    =
    \int_{t_{\mathrm{emit}}}^{t_{\mathrm{obs}}}
    \Gamma_{\mathrm{D}}(t)
    \,\mathrm{d}t.
    \label{eq:photon_relaxation_exact_time_redshift}
\end{equation}

Let \(\ell\) denote arc length along the physical ray and define
\begin{equation}
    v_{\mathrm{ray}}(\ell)
    =
    \left|
        \mathbf{v}_{\mathrm{g}}(\ell)
    \right|.
    \label{eq:photon_relaxation_ray_speed}
\end{equation}
The propagation time is
\begin{equation}
    t_{\mathrm{obs}}
    -
    t_{\mathrm{emit}}
    =
    \int_{\mathrm{path}}
    \frac{
        \mathrm{d}\ell
    }{
        v_{\mathrm{ray}}
        \left[
            \omega(\ell),
            \alpha_{\mathrm{G}}(\ell)
        \right]
    }.
    \label{eq:photon_relaxation_path_time}
\end{equation}
Therefore, for a stationary but spatially nonuniform path,
\begin{equation}
    \ln
    \left(
        1+z
    \right)
    =
    \int_{\mathrm{path}}
    \frac{
        \Gamma_{\mathrm{D}}(\ell)
    }{
        v_{\mathrm{ray}}
        \left[
            \omega(\ell),
            \alpha_{\mathrm{G}}(\ell)
        \right]
    }
    \,\mathrm{d}\ell.
    \label{eq:photon_relaxation_general_path_redshift}
\end{equation}
This expression separates the conservative propagation time from the nonconservative relaxation accumulated per unit time.

In the weakly loaded constant-rate limit,
\begin{equation}
    \alpha_{\mathrm{G}}
    \rightarrow
    0,
    \qquad
    v_{\mathrm{ray}}
    \simeq
    c_{\mathrm{E}},
    \qquad
    \Gamma_{\mathrm{D}}
    \simeq
    \text{constant},
    \label{eq:photon_relaxation_void_limit}
\end{equation}
Equation~\eqref{eq:photon_relaxation_general_path_redshift} reduces to
\begin{equation}
    1+z
    =
    \exp
    \left(
        \frac{
            \Gamma_{\mathrm{D}}r
        }{
            c_{\mathrm{E}}
        }
    \right).
    \label{eq:photon_relaxation_exact_distance_redshift}
\end{equation}
Define the characteristic relaxation length
\begin{equation}
    R_{\mathrm{U}}
    =
    \frac{
        c_{\mathrm{E}}
    }{
        \Gamma_{\mathrm{D}}
    }.
    \label{eq:photon_relaxation_ru_definition}
\end{equation}
Then
\begin{equation}
    1+z
    =
    e^{r/R_{\mathrm{U}}},
    \qquad
    r
    =
    R_{\mathrm{U}}
    \ln
    \left(
        1+z
    \right).
    \label{eq:photon_relaxation_exact_euclidean_distance}
\end{equation}

The low-distance expansion is
\begin{equation}
    z
    =
    \frac{
        r
    }{
        R_{\mathrm{U}}
    }
    +
    \frac{1}{2}
    \left(
        \frac{
            r
        }{
            R_{\mathrm{U}}
        }
    \right)^2
    +
    \mathcal{O}
    \left(
        \frac{
            r^3
        }{
            R_{\mathrm{U}}^3
        }
    \right),
    \label{eq:photon_relaxation_low_redshift_expansion}
\end{equation}
so
\begin{equation}
    c_{\mathrm{E}}z
    \simeq
    \Gamma_{\mathrm{D}}r
    \label{eq:photon_relaxation_linear_redshift_limit}
\end{equation}
is only the leading low-redshift approximation.

The absolute magnitude of \(\Gamma_{\mathrm{D}}\) is not yet derived from the FCC--Condensate microphysics. At the present level of TFT, it is fixed empirically by matching Equation~\eqref{eq:photon_relaxation_linear_redshift_limit} to the observed local Hubble slope,
\begin{equation}
    c_{\mathrm{E}}z
    \simeq
    H_{0}^{(\mathrm{obs})}r.
\end{equation}
Therefore, within the TFT low-redshift calibration,
\begin{equation}
    \Gamma_{\mathrm{D}}^{(\mathrm{cal})}
    =
    H_{0}^{(\mathrm{obs})}.
    \label{eq:photon_relaxation_gamma_calibration}
\end{equation}
Using the adopted local value
\begin{equation}
    H_{0}^{(\mathrm{obs})}
    =
    73.04\,
    \mathrm{km\,s^{-1}\,Mpc^{-1}},
\end{equation}
gives
\begin{equation}
    \Gamma_{\mathrm{D}}^{(\mathrm{cal})}
    \approx
    2.3671
    \times
    10^{-18}
    \ \mathrm{s}^{-1}.
    \label{eq:photon_relaxation_gamma_numerical}
\end{equation}
Thus, the equality between \(\Gamma_{\mathrm{D}}\) and the observed low-redshift Hubble slope is an empirical normalization of the TFT relaxation coefficient, not an independent microscopic derivation. A first-principles determination of \(\Gamma_{\mathrm{D}}\) from the causal FCC--Condensate response remains a theory-completion requirement.

Preservation of coherent modal action also gives
\begin{equation}
    \frac{
        E_{\mathrm{obs}}
    }{
        E_{\mathrm{emit}}
    }
    =
    \frac{
        \omega_{\mathrm{obs}}
    }{
        \omega_{\mathrm{emit}}
    }
    =
    \frac{1}{1+z}.
    \label{eq:photon_relaxation_energy_redshift}
\end{equation}
This is the energy reduction of the surviving coherent carrier under the selected transport closure. It must not be supplemented by an independent amplitude-extinction factor unless a separate absorption or incoherent-loss channel is introduced.

Equation~\eqref{eq:photon_relaxation_exact_euclidean_distance} is exact within the constant-\(\Gamma_{\mathrm{D}}\), weak-loading, stationary Euclidean-path closure. The general result remains Equation~\eqref{eq:photon_relaxation_general_path_redshift}. Moreover, \(R_{\mathrm{U}}\) is an \(e\)-folding length, not a universal stopping radius or zero-energy boundary.


\section{Observational Screening and Self-Consistent Path Integration}
\label{sec:photon_relaxation_observational_screening}

A viable photon branch must remain compatible with the absence of a resolved energy-ordered delay in high-energy transient observations. Two related tests are applied below: a local comparison with the conventional quadratic time-of-flight parametrization and a TFT path integral that includes the frequency evolution along the propagation path.

In a weakly loaded region,
\begin{equation}
    E
    \simeq
    \hbar c_{\mathrm{E}}k.
    \label{eq:photon_relaxation_energy_wave_number}
\end{equation}
Using
\begin{equation}
    d_{\mathrm{nn}}
    =
    2r_{\mathrm{D}},
    \qquad
    E_{\max}
    =
    \frac{
        hc_{\mathrm{E}}
    }{
        2r_{\mathrm{D}}
    }
    =
    \frac{
        hc_{\mathrm{E}}
    }{
        d_{\mathrm{nn}}
    },
    \label{eq:photon_relaxation_energy_scale}
\end{equation}
one obtains
\begin{equation}
    kd_{\mathrm{nn}}
    =
    2\pi
    \frac{
        E
    }{
        E_{\max}
    }.
    \label{eq:photon_relaxation_dimensionless_energy}
\end{equation}
Equivalently, the natural low-wave-number conversion scale is
\[
    E_d
    \equiv
    \frac{
        \hbar c_{\mathrm{E}}
    }{
        d_{\mathrm{nn}}
    }
    =
    \frac{
        E_{\max}
    }{
        2\pi
    },
    \qquad
    kd_{\mathrm{nn}}
    \simeq
    \frac{E}{E_d}.
\]
The notation \(E_{\max}\) is retained because it is the established TFT calibration scale. In this calculation it sets the wavelength normalization; it is not identified with the exact energy at a Bloch-zone boundary.

For the baseline nearest-neighbor branch, Equation~\eqref{eq:photon_relaxation_baseline_group_velocity_expansion} gives
\begin{equation}
    1
    -
    \frac{
        v_{\mathrm{g}}
    }{
        c_{\mathrm{P}}
    }
    =
    \frac{
        \pi^2
    }{
        2
    }
    \left(
        \frac{
            E
        }{
            E_{\max}
        }
    \right)^2
    +
    \mathcal{O}
    \left(
        \frac{
            E^4
        }{
            E_{\max}^4
        }
    \right).
    \label{eq:photon_relaxation_baseline_energy_delay}
\end{equation}

The conventional subluminal quadratic comparison law is
\begin{equation}
    \frac{
        v(E)
    }{
        c_{\mathrm{E}}
    }
    \simeq
    1
    -
    \frac{3}{2}
    \left(
        \frac{
            E
        }{
            E_{\mathrm{QG},2}
        }
    \right)^2,
    \label{eq:photon_relaxation_standard_quadratic_parameterization}
\end{equation}
where \(E_{\mathrm{QG},2}\) is an observational comparison scale rather than a TFT material constant. In the weak-loading limit
\(c_{\mathrm{P}}\simeq c_{\mathrm{E}}\), matching the quadratic coefficients gives
\begin{equation}
    E_{\mathrm{QG},2}^{(\mathrm{base})}
    =
    \frac{
        \sqrt{3}
    }{
        \pi
    }
    E_{\max}.
    \label{eq:photon_relaxation_baseline_effective_scale}
\end{equation}

For the adopted TFT calibration
\begin{equation}
    E_{\max}
    \approx
    8.320
    \times
    10^{10}
    \ \mathrm{GeV},
    \label{eq:photon_relaxation_maximum_energy_numerical}
\end{equation}
the baseline branch corresponds to
\begin{equation}
    E_{\mathrm{QG},2}^{(\mathrm{base})}
    \approx
    4.59
    \times
    10^{10}
    \ \mathrm{GeV}.
    \label{eq:photon_relaxation_baseline_effective_scale_numerical}
\end{equation}

The LHAASO analysis of GRB~221009A reports
\begin{equation}
    E_{\mathrm{QG},2}
    >
    6
    \times
    10^{-8}
    E_{\mathrm{Pl}}
    \approx
    7.3
    \times
    10^{11}
    \ \mathrm{GeV}
    \qquad
    \text{at \(95\%\) confidence}
    \label{eq:photon_relaxation_lhaaso_quadratic_bound}
\end{equation}
for its quadratic Lorentz-invariance-violation parametrization
\cite{LHAASO2024LIV}. The local coefficient comparison therefore places the baseline nearest-neighbor photon branch in strong tension with the selected observational bound.

A self-consistent TFT test must also account for the changing carrier energy along the path. Consider a weakly loaded path of length \(r\), parameterized by \(\ell\), with the observer at \(\ell=r\). Under the constant-rate transport law,
\begin{equation}
    E(\ell)
    =
    E_{\mathrm{obs}}
    \exp
    \left(
        \frac{
            r-\ell
        }{
            R_{\mathrm{U}}
        }
    \right),
    \label{eq:photon_relaxation_path_energy}
\end{equation}
and
\begin{equation}
    1+z_{\mathrm{s}}
    =
    e^{r/R_{\mathrm{U}}}.
    \label{eq:photon_relaxation_path_scale}
\end{equation}

Suppose the leading local fractional speed reduction is
\begin{equation}
    1
    -
    \frac{
        v_{\mathrm{g}}(E)
    }{
        c_{\mathrm{E}}
    }
    =
    A_n
    \left(
        \frac{
            E
        }{
            E_{\max}
        }
    \right)^n.
    \label{eq:photon_relaxation_generic_fractional_delay}
\end{equation}
The arrival-time difference between simultaneously emitted components with observed energies
\(E_{\mathrm{h,obs}}\) and \(E_{\mathrm{l,obs}}\) is then
\begin{align}
    \Delta t_n
    &\simeq
    \frac{1}{c_{\mathrm{E}}}
    \int_0^r
    \left[
        \delta_n
        \left(
            E_{\mathrm{h}}(\ell)
        \right)
        -
        \delta_n
        \left(
            E_{\mathrm{l}}(\ell)
        \right)
    \right]
    \mathrm{d}\ell
    \nonumber\\
    &=
    \frac{
        A_n
    }{
        n\Gamma_{\mathrm{D}}
    }
    \left[
        \left(
            \frac{
                E_{\mathrm{h,obs}}
            }{
                E_{\max}
            }
        \right)^n
        -
        \left(
            \frac{
                E_{\mathrm{l,obs}}
            }{
                E_{\max}
            }
        \right)^n
    \right]
    \left[
        \left(
            1+z_{\mathrm{s}}
        \right)^n
        -
        1
    \right],
    \label{eq:photon_relaxation_integrated_delay_general}
\end{align}
where
\(\delta_n(E)=A_n(E/E_{\max})^n\).
Equation~\eqref{eq:photon_relaxation_integrated_delay_general} is conditional on the same weak-loading, constant-\(\Gamma_{\mathrm{D}}\) transport closure as Equation~\eqref{eq:photon_relaxation_exact_euclidean_distance}.

For the baseline branch,
\begin{equation}
    n
    =
    2,
    \qquad
    A_2
    =
    \frac{
        \pi^2
    }{
        2
    }.
    \label{eq:photon_relaxation_baseline_delay_coefficient}
\end{equation}
For GRB~221009A, use
\begin{equation}
    z_{\mathrm{s}}
    =
    0.151,
    \qquad
    E_{\mathrm{l,obs}}
    =
    0.2
    \ \mathrm{TeV},
    \qquad
    E_{\mathrm{h,obs}}
    =
    7
    \ \mathrm{TeV},
    \label{eq:photon_relaxation_lhaaso_screening_inputs}
\end{equation}
together with the empirical TFT calibration in Equation~\eqref{eq:photon_relaxation_gamma_numerical}. Equation~\eqref{eq:photon_relaxation_integrated_delay_general} then gives
\begin{equation}
    \Delta t_2^{(\mathrm{base})}
    \approx
    2.4
    \times
    10^3
    \ \mathrm{s}.
    \label{eq:photon_relaxation_baseline_integrated_delay}
\end{equation}
This macroscopically large delay strongly disfavors identifying the reduced baseline branch with the completed physical photon branch under the selected TFT propagation law.

For the dispersion-suppressed branch, Equations~\eqref{eq:photon_relaxation_suppressed_group_velocity} and~\eqref{eq:photon_relaxation_fcc_directional_bounds} give
\begin{equation}
    1
    -
    \frac{
        v_{\mathrm{g},\parallel}
    }{
        c_{\mathrm{P}}
    }
    \leq
    \frac{
        17
    }{
        1152
    }
    \left(
        kd_{\mathrm{nn}}
    \right)^4
    +
    \mathcal{O}
    \left(
        k^6d_{\mathrm{nn}}^6
    \right).
    \label{eq:photon_relaxation_suppressed_upper_bound}
\end{equation}
In energy form,
\begin{equation}
    1
    -
    \frac{
        v_{\mathrm{g},\parallel}
    }{
        c_{\mathrm{P}}
    }
    \leq
    \frac{
        17\pi^4
    }{
        72
    }
    \left(
        \frac{
            E
        }{
            E_{\max}
        }
    \right)^4
    +
    \mathcal{O}
    \left(
        \frac{
            E^6
        }{
            E_{\max}^6
        }
    \right).
    \label{eq:photon_relaxation_suppressed_energy_bound}
\end{equation}
Thus,
\begin{equation}
    n
    =
    4,
    \qquad
    A_4
    =
    \frac{
        17\pi^4
    }{
        72
    }.
    \label{eq:photon_relaxation_suppressed_delay_coefficient}
\end{equation}
Using the same source redshift, observed-energy interval, and relaxation calibration gives
\begin{equation}
    \Delta t_4^{(\mathrm{sup})}
    \lesssim
    9.2
    \times
    10^{-11}
    \ \mathrm{s}.
    \label{eq:photon_relaxation_suppressed_integrated_delay}
\end{equation}
The selected dispersion-suppressed branch therefore produces a negligible path-integrated chromatic delay over the tested TeV interval.

This is a kinematic screening calculation, not a detector-level likelihood analysis. A complete observational test must include the measured source light curve, intrinsic spectral evolution, detector response, and energy uncertainty. Within the present screening scope, the baseline nearest-neighbor branch is strongly disfavored, whereas the selected quartic branch remains compatible with the quoted timing benchmark. The comparison constrains the effective photon eigenbranch but does not derive the microscopic origin of the \(1/12\) coefficient.

\section{Finite-Packet Cycle Dilation and Its Scope}
\label{sec:photon_relaxation_packet_dilation_scope}

Preservation of the coherent modal action
\(E/\omega\)
does not by itself determine the temporal width of a finite packet. Packet-duration evolution requires an additional condition on the packet phase structure.

For a packet bounded by \(t_{-}\) and \(t_{+}\), define the carrier-cycle count by
\begin{equation}
    N
    =
    \int_{t_{-}}^{t_{+}}
    f(t)
    \,\mathrm{d}t.
    \label{eq:photon_relaxation_packet_cycle_count}
\end{equation}
For a narrowband packet whose carrier frequency varies negligibly across its own support,
\begin{equation}
    N
    \simeq
    f_{\mathrm{c}}
    \Delta t_{\mathrm{pkt}},
    \qquad
    \Delta t_{\mathrm{pkt}}
    =
    t_{+}
    -
    t_{-}.
    \label{eq:photon_relaxation_narrowband_cycle_count}
\end{equation}

The additional fixed-cycle closure is
\begin{equation}
    N_{\mathrm{obs}}
    =
    N_{\mathrm{emit}}.
    \label{eq:photon_relaxation_packet_cycle_conservation}
\end{equation}
This condition is distinct from modal-action preservation. It states that the packet retains its ordered carrier-phase content during propagation.

If the frequency relaxation is effectively uniform across the packet,
\begin{equation}
    f_{\mathrm{emit}}
    \Delta t_{\mathrm{pkt,emit}}
    =
    f_{\mathrm{obs}}
    \Delta t_{\mathrm{pkt,obs}}.
    \label{eq:photon_relaxation_packet_cycle_balance}
\end{equation}
Using
\begin{equation}
    \frac{
        f_{\mathrm{emit}}
    }{
        f_{\mathrm{obs}}
    }
    =
    1+z,
    \label{eq:photon_relaxation_frequency_redshift}
\end{equation}
gives
\begin{equation}
    \Delta t_{\mathrm{pkt,obs}}
    =
    \left(
        1+z
    \right)
    \Delta t_{\mathrm{pkt,emit}}.
    \label{eq:photon_relaxation_packet_duration_dilation}
\end{equation}

Equation~\eqref{eq:photon_relaxation_packet_duration_dilation} therefore requires preservation of cycle count, narrowband or adiabatic carrier evolution, approximately uniform relaxation across the packet, and survival of the packet as one coherent phase structure. It does not establish dilation of an arbitrary macroscopic source light curve.

To show the distinction, consider two independently emitted packets following the same stationary path:
\begin{equation}
    t_{\mathrm{obs}}^{(j)}
    =
    t_{\mathrm{emit}}^{(j)}
    +
    t_{\mathrm{prop}}^{(j)},
    \qquad
    j
    =
    1,2.
    \label{eq:photon_relaxation_two_packet_arrivals}
\end{equation}
Their observed separation is
\begin{equation}
    \Delta t_{\mathrm{obs}}
    =
    \Delta t_{\mathrm{emit}}
    +
    t_{\mathrm{prop}}^{(2)}
    -
    t_{\mathrm{prop}}^{(1)}.
    \label{eq:photon_relaxation_arrival_interval_general}
\end{equation}
For identical spectral components emitted along the same stationary path,
\begin{equation}
    t_{\mathrm{prop}}^{(2)}
    =
    t_{\mathrm{prop}}^{(1)},
    \label{eq:photon_relaxation_equal_stationary_path_times}
\end{equation}
so
\begin{equation}
    \Delta t_{\mathrm{obs}}
    =
    \Delta t_{\mathrm{emit}}.
    \label{eq:photon_relaxation_no_automatic_envelope_dilation}
\end{equation}

Carrier-frequency relaxation and finite-packet dilation therefore do not automatically imply
\begin{equation}
    \Delta t_{\mathrm{light\,curve,obs}}
    =
    \left(
        1+z
    \right)
    \Delta t_{\mathrm{light\,curve,emit}}
    \label{eq:photon_relaxation_macroscopic_dilation_target}
\end{equation}
for a sequence of independently emitted packets.

A derivation of astronomical transient dilation requires an additional source-to-envelope transport law. That law must either show that the macroscopic emission envelope is itself carried by a coherent finite-cycle structure or derive an equivalent mapping between source chronology and observer chronology. The residual lattice dispersion evaluated in Section~\ref{sec:photon_relaxation_observational_screening} is far too small to produce an order-\((1+z)\) stretching of a broadband transient envelope.

The result established here is therefore restricted to
\begin{equation}
    \begin{aligned}
        &\text{fixed-cycle narrowband coherent packet}
        \\
        &\hspace{5em}
        \longrightarrow
        \text{conditional packet-duration dilation}.
    \end{aligned}
    \label{eq:photon_relaxation_packet_dilation_scope_result}
\end{equation}
Macroscopic light-curve dilation remains a separate cosmological transport-completion problem.

\section{Threshold Assimilation into a CMB-Associated Thermal Reservoir}
\label{sec:photon_relaxation_cmb_assimilation}

The exponential transport law does not drive the coherent carrier to zero frequency or energy at any finite distance. In the weak-loading, constant-rate limit,
\begin{equation}
    \omega(r)
    =
    \omega_{\mathrm{emit}}
    \exp
    \left(
        -
        \frac{r}{R_{\mathrm{U}}}
    \right),
    \label{eq:photon_relaxation_frequency_distance_law}
\end{equation}
and, because
\(\mathcal{I}_{\gamma}=E/\omega\)
is preserved,
\begin{equation}
    E(r)
    =
    E_{\mathrm{emit}}
    \exp
    \left(
        -
        \frac{r}{R_{\mathrm{U}}}
    \right).
    \label{eq:photon_relaxation_energy_distance_law}
\end{equation}
Consequently,
\begin{equation}
    E(R_{\mathrm{U}})
    =
    \frac{
        E_{\mathrm{emit}}
    }{
        e
    }.
    \label{eq:photon_relaxation_one_efold_energy}
\end{equation}
Thus, \(R_{\mathrm{U}}\) is an \(e\)-folding length, not a universal zero-energy or stopping radius.

A finite coherence horizon may nevertheless arise if the relaxed transverse mode becomes dynamically indistinguishable from the ambient thermal spectrum. Because the CMB is broadband, introduce an effective coherence-loss frequency
\begin{equation}
    \omega_{\mathrm{CMB}}^{(\mathrm{eff})}
    =
    2\pi
    f_{\mathrm{CMB}}^{(\mathrm{eff})}.
    \label{eq:photon_relaxation_effective_cmb_frequency}
\end{equation}
This parameter represents an effective spectral and coherence boundary, not a unique frequency of the CMB.

For a carrier with modal action
\(\mathcal{I}_{\gamma}\),
the remaining threshold energy is
\begin{equation}
    E_{\mathrm{th}}^{(\mathrm{eff})}
    =
    \mathcal{I}_{\gamma}
    \omega_{\mathrm{CMB}}^{(\mathrm{eff})}.
    \label{eq:photon_relaxation_effective_cmb_energy}
\end{equation}
The threshold energy therefore depends on the action carried by the coherent mode. For one quantum-normalized excitation,
\begin{equation}
    \mathcal{I}_{\gamma}
    =
    \hbar,
    \label{eq:photon_relaxation_quantum_normalized_action}
\end{equation}
and Equation~\eqref{eq:photon_relaxation_effective_cmb_energy} reduces to
\[
    E_{\mathrm{th}}^{(\mathrm{eff})}
    =
    \hbar
    \omega_{\mathrm{CMB}}^{(\mathrm{eff})}
    =
    h
    f_{\mathrm{CMB}}^{(\mathrm{eff})}.
\]

The assimilation condition is most fundamentally stated in frequency:
\begin{equation}
    \omega
    \left(
        r_{\mathrm{assim}}
    \right)
    =
    \omega_{\mathrm{CMB}}^{(\mathrm{eff})},
    \label{eq:photon_relaxation_assimilation_condition}
\end{equation}
with
\begin{equation}
    \omega_{\mathrm{emit}}
    >
    \omega_{\mathrm{CMB}}^{(\mathrm{eff})}.
    \label{eq:photon_relaxation_assimilation_domain}
\end{equation}
Using Equation~\eqref{eq:photon_relaxation_frequency_distance_law},
\begin{equation}
    r_{\mathrm{assim}}
    =
    R_{\mathrm{U}}
    \ln
    \left(
        \frac{
            \omega_{\mathrm{emit}}
        }{
            \omega_{\mathrm{CMB}}^{(\mathrm{eff})}
        }
    \right)
    =
    R_{\mathrm{U}}
    \ln
    \left(
        \frac{
            f_{\mathrm{emit}}
        }{
            f_{\mathrm{CMB}}^{(\mathrm{eff})}
        }
    \right).
    \label{eq:photon_relaxation_assimilation_distance_frequency}
\end{equation}
Because the same modal action multiplies the initial and threshold frequencies,
\begin{equation}
    r_{\mathrm{assim}}
    =
    R_{\mathrm{U}}
    \ln
    \left(
        \frac{
            E_{\mathrm{emit}}
        }{
            E_{\mathrm{th}}^{(\mathrm{eff})}
        }
    \right).
    \label{eq:photon_relaxation_assimilation_distance_energy}
\end{equation}

The assimilation distance is therefore frequency dependent. For
\begin{equation}
    \frac{
        f_{\mathrm{emit}}
    }{
        f_{\mathrm{CMB}}^{(\mathrm{eff})}
    }
    \sim
    10^3
    \text{--}
    10^4,
    \label{eq:photon_relaxation_optical_cmb_ratio}
\end{equation}
one obtains
\begin{equation}
    r_{\mathrm{assim}}
    \sim
    \left(
        7
        \text{--}
        9
    \right)
    R_{\mathrm{U}}.
    \label{eq:photon_relaxation_optical_assimilation_scale}
\end{equation}
For the adopted TFT calibration of \(R_{\mathrm{U}}\), this is of order
\begin{equation}
    10^{27}
    \ \mathrm{m}.
    \label{eq:photon_relaxation_assimilation_order}
\end{equation}
This is an approximate optical-to-thermal coherence scale, not a universal photon horizon.

\subsection{Energy Accounting at the Assimilation Threshold}
\label{subsec:photon_relaxation_assimilation_energy_accounting}

Before terminal coherence loss, integration of the carrier--reservoir balance gives
\begin{equation}
    E_{\mathrm{res}}(t)
    -
    E_{\mathrm{res}}
    \left(
        t_{\mathrm{emit}}
    \right)
    =
    E_{\mathrm{emit}}
    -
    E(t),
    \label{eq:photon_relaxation_integrated_reservoir_energy}
\end{equation}
provided that no additional absorptive, radiative, or mode-conversion channel is active.

Immediately before assimilation,
\begin{equation}
    E
    \left(
        r_{\mathrm{assim}}^{-}
    \right)
    =
    E_{\mathrm{th}}^{(\mathrm{eff})}.
    \label{eq:photon_relaxation_preassimilation_energy}
\end{equation}
The continuously transferred energy is therefore
\begin{equation}
    \Delta E_{\mathrm{res}}^{(\mathrm{cont})}
    =
    E_{\mathrm{emit}}
    -
    E_{\mathrm{th}}^{(\mathrm{eff})}.
    \label{eq:photon_relaxation_continuous_reservoir_transfer}
\end{equation}
Under the terminal assimilation closure, the remaining coherent energy is redistributed into accessible reservoir modes:
\begin{equation}
    \Delta E_{\mathrm{res}}^{(\mathrm{assim})}
    =
    E_{\mathrm{th}}^{(\mathrm{eff})}.
    \label{eq:photon_relaxation_terminal_assimilation_energy}
\end{equation}
Hence
\begin{equation}
    \Delta E_{\mathrm{res}}^{(\mathrm{tot})}
    =
    \Delta E_{\mathrm{res}}^{(\mathrm{cont})}
    +
    \Delta E_{\mathrm{res}}^{(\mathrm{assim})}
    =
    E_{\mathrm{emit}}.
    \label{eq:photon_relaxation_total_reservoir_transfer}
\end{equation}

Equation~\eqref{eq:photon_relaxation_total_reservoir_transfer} closes only the energy bookkeeping of one assimilated coherent mode. It does not establish global thermal equilibrium, detailed balance, or preservation of the measured CMB spectrum. Those require a microscopic theory of the terminal instability, the destination of deposited energy and momentum, and the resulting thermal and spectral response.

\section{Closure Status, Falsifiability, and Remaining Completion}
\label{sec:photon_relaxation_status}

The appendix contains three logically distinct elements: results derived from the selected reduced carrier action, conditional nonconservative transport closures, and unresolved microscopic or cosmological completion requirements.

\subsection{Results Derived Within the Selected Reduced Carrier Action}
\label{subsec:photon_relaxation_derived_status}

Within the divergence-free positive two-term action of Section~\ref{sec:photon_relaxation_minimal_closure}, the following results are derived:

\begin{enumerate}
    \item The baseline directional chain has a leading group-speed correction proportional to
    \(k^2d_{\mathrm{nn}}^2\).

    \item The positive centered-second-difference term with coefficient \(1/12\) cancels the quartic term in the expansion of \(\omega^2\) within the selected stencil.

    \item The resulting reduced transverse branch is gapless, inversion symmetric, spectrally nonnegative, and degenerate in its two transverse components.

    \item Its leading radial group-speed correction is proportional to
    \(k^4d_{\mathrm{nn}}^4\), with the FCC directional bound in Equation~\eqref{eq:photon_relaxation_fcc_directional_bounds}.

    \item The GRB screening calculation strongly disfavors the baseline nearest-neighbor branch while leaving the selected quartic branch compatible with the adopted timing benchmark.
\end{enumerate}

These results follow from the selected reduced action; they do not establish that the complete FCC parent dynamics uniquely generate that action.

\subsection{Conditional Nonconservative Closures}
\label{subsec:photon_relaxation_conditional_status}

The relaxation-redshift construction additionally assumes:

\begin{enumerate}
    \item a causal weak-loss response with the inverse-frequency loss tangent in Equation~\eqref{eq:photon_relaxation_inverse_frequency_phase_lag};

    \item preservation of the coherent modal action
    \(\mathcal{I}_{\gamma}=E/\omega\);

    \item the direction-preserving wave-vector transport law in Equation~\eqref{eq:photon_relaxation_longitudinal_k_drag};

    \item canonical energy--quasimomentum balance with the FCC--Condensate reservoir;

    \item suppression of incoherent absorption, angular diffusion, phase noise, polarization scrambling, and secondary-mode conversion;

    \item a stationary or quasistatic background, with approximately constant
    \(\Gamma_{\mathrm{D}}\)
    for the simple Euclidean exponential law;

    \item fixed carrier-cycle count for conditional packet-duration dilation; and

    \item an effective thermal coherence threshold for terminal assimilation.
\end{enumerate}

Under these conditions, modal-energy transfer produces coherent frequency evolution. The simple distance law follows only after the additional weak-loading and constant-rate approximations are imposed.

\subsection{Results Not Established by the Present Appendix}
\label{subsec:photon_relaxation_not_established}

The present construction does not establish:

\begin{enumerate}
    \item unique microscopic origin of the coefficient \(1/12\);

    \item uniform finite-\(k\) endpoint loading from the fully resolved FCC dynamic matrix;

    \item first-principles derivation of the inverse-frequency loss tangent;

    \item nonlinear proof of coherent modal-action preservation;

    \item derivation of the wave-vector transport law from the same causal susceptibility that produces the phase lag;

    \item quantitative conversion of canonical Bloch quasimomentum into mechanical substrate momentum;

    \item suppression of cumulative broadening, diffusion, polarization change, recoil radiation, and secondary modes;

    \item dilation of arbitrary macroscopic transient light curves;

    \item derivation of
    \(f_{\mathrm{CMB}}^{(\mathrm{eff})}\)
    from the measured thermal spectrum; or

    \item global thermal, spectral, and momentum balance of the deposited radiation.
\end{enumerate}

\subsection{Falsifiable Consequences}
\label{subsec:photon_relaxation_falsifiability}

The construction is testable through the following channels:

\begin{enumerate}
    \item \textbf{Photon dispersion.}
    A confirmed leading quadratic vacuum delay at the TFT lattice scale would contradict the selected dispersion-suppressed branch.

    \item \textbf{FCC directional structure.}
    The leading residual correction carries the cubic factor in Equation~\eqref{eq:photon_relaxation_fcc_directional_bounds}; its interpretation depends on the orientation or domain structure of the vacuum lattice.

    \item \textbf{Loss-tangent scaling.}
    A frequency-independent
    \(\Gamma_{\mathrm{D}}\)
    requires the inverse-frequency weak-loss behavior derived above.

    \item \textbf{Spectral and directional coherence.}
    Excessive line broadening, frequency diffusion, image blurring, polarization randomization, or secondary radiation would contradict the assumed coherent transport channel.

    \item \textbf{Momentum balance.}
    Failure to construct a mechanical substrate response consistent with the required canonical quasimomentum change would invalidate the wave-vector transport closure.

    \item \textbf{Macroscopic time dilation.}
    Failure to derive source-envelope dilation elsewhere in TFT would leave astronomical transient-duration observations unexplained.

    \item \textbf{Assimilation scale.}
    The selected closure predicts a frequency-dependent coherence distance rather than a universal stopping radius.

    \item \textbf{Thermal consistency.}
    Unacceptable CMB spectral distortion or secular reservoir heating would contradict the proposed energy-deposition mechanism.
\end{enumerate}

\subsection{Required Microscopic and Cosmological Completion}
\label{subsec:photon_relaxation_remaining_completion}

The remaining program is to:

\begin{enumerate}
    \item derive the physical finite-wave-number transverse branch from the complete FCC contact Hessian, rotational variables, internal phases, and Tension-Column coupling;

    \item determine the loaded finite-\(k\) dynamic matrix without importing the long-wavelength endpoint-transfer factor uniformly;

    \item construct a causal FCC--Condensate susceptibility satisfying the appropriate Kramers--Kronig and fluctuation--dissipation constraints;

    \item derive the real-frequency and real-wave-vector transport law from that same response;

    \item derive the conversion between Bloch quasimomentum and mechanical substrate momentum;

    \item demonstrate nonlinear stability against absorption, packet rupture, diffusion, decoherence, polarization scrambling, and secondary-mode production;

    \item determine
    \(\Gamma_{\mathrm{D}}\)
    independently from vacuum microphysics;

    \item derive macroscopic source-envelope transport and test astronomical light-curve dilation;

    \item derive
    \(f_{\mathrm{CMB}}^{(\mathrm{eff})}\)
    and the instability responsible for terminal coherence loss; and

    \item establish global thermal, spectral, energy, and momentum balance for the deposited radiation.
\end{enumerate}

The logical hierarchy is
\begin{align}
    &\text{complete FCC contact and internal-phase dynamics}
    \nonumber\\
    &\qquad\longrightarrow
    \text{physical finite-\(k\) transverse eigenbranch}
    \nonumber\\
    &\qquad\longrightarrow
    \text{causal FCC--Condensate response}
    \nonumber\\
    &\qquad\longrightarrow
    \text{coherent action and wave-vector transport}
    \nonumber\\
    &\qquad\longrightarrow
    \text{relaxation redshift and threshold assimilation}.
    \label{eq:photon_relaxation_final_hierarchy}
\end{align}

The appendix therefore supplies a dimensionally closed reduced transverse carrier action and an explicit conditional nonconservative transport closure while maintaining a strict separation among inherited TFT structure, derived results, constitutive assumptions, observational consequences, and unresolved completion.